% ============================================================================
\section{Introduction: Why Rethink Neural Computation?}
\label{sec:introduction}
% ============================================================================

Every time you ask a voice assistant a question, every time your phone
recognises a face, every time a language model generates a sentence,
trillions of tiny multiplications happen inside specialised hardware.
These multiplications---floating-point numbers multiplied by
floating-point weights and summed up---are the heartbeat of modern AI.
The industry calls this operation \concept{multiply-accumulate} (MAC),
and virtually every neural network alive today is built on it.

This dependence has consequences---and the industry knows it.
\emph{Quantisation} has become the standard countermeasure: reducing
weights from 32-bit floating-point to 8-bit or even 4-bit integers
cuts memory and energy dramatically. A 70-billion-parameter model
that once required 140\,GB in 16-bit format now fits in roughly
35\,GB at 4~bits (via methods such as GPTQ~\cite{frantar2022gptq} or
AWQ~\cite{lin2023awq}), bringing it within reach of consumer
hardware. Experimental 2-bit quantisation pushes this below 18\,GB.

Yet even aggressively quantised models still rely on integer
multiplication and addition---operations that, while far cheaper than
floating-point, remain significantly more expensive than the simplest
thing a processor can do:

\begin{itemize}[leftmargin=2em]
  \item \textbf{Energy:} At 45\,nm, a 32-bit floating-point
    multiplication costs 3.7\,pJ; an 8-bit integer multiplication costs
    0.2\,pJ---an $18\times$ reduction. But a single bitwise logic
    operation (AND, OR, XOR) costs on the order of 0.1\,pJ, cheaper
    still and $37\times$ below floating-point~\cite{horowitz2014energy}.
    These ratios reflect the algorithmic complexity of each operation
    (a multiplier requires $O(n^2)$ gates, an adder $O(n)$, a bitwise
    operation $O(1)$ per bit) and therefore hold across process
    nodes~\cite{stillmaker2017scaling}.
  \item \textbf{Hardware monopoly:} Training large models still requires
    GPUs, a market dominated by a single manufacturer. Quantised
    inference can run on CPUs, but the underlying matrix multiplications
    over integer values still benefit heavily from specialised
    accelerators.
  \item \textbf{Deployment barriers:} Even with 4-bit quantisation, a
    70-billion-parameter model requires roughly 35\,GB---feasible on a
    high-end laptop but not on a microcontroller, a sensor, or a phone
    without dedicated AI silicon. A fully binary network storing one bit
    per weight would need only $\sim$8\,GB for the same model.
\end{itemize}

\begin{keyinsight}
The question this article explores is not ``How do we build faster
GPUs?'' but rather ``Can we build neural networks that do not need
GPUs at all?''---networks that run on the bitwise logic instructions
every CPU already has.
\end{keyinsight}

\begin{engineernote}
This is a \emph{CPU-first} question, not a claim that a Boolean circuit
automatically runs quickly on a CPU. A useful building block needs packed,
SIMD-friendly data layout, predictable control flow, and local memory
access. The relevant comparison is end-to-end latency and energy against a
small quantised integer baseline on the same processor.
\end{engineernote}

\subsection{The Four Paradigms}

Over the past decade, four distinct research directions have emerged
that aim to replace MAC operations with \emph{bitwise} neural
computation---each attacking the problem from a different angle.
\Cref{fig:paradigm-map} shows how they relate.

\begin{figure}[htbp]
\centering
\begin{tikzpicture}[
  box/.style={
    draw, rounded corners=4pt, minimum width=3.8cm, minimum height=1.6cm,
    align=center, font=\small\sffamily, line width=0.8pt
  },
  arrow/.style={-{Stealth[length=5pt]}, thick, bitgray},
  label/.style={font=\scriptsize\sffamily, bitgray}
]
  % Boxes -- wider spacing for readable arrow labels
  \node[box, fill=bitblue!10, draw=bitblue] (bnn)
    at (0,0) {\textbf{Binary Neural}\\\textbf{Networks}\\[2pt]\scriptsize XNOR + popcount};
  \node[box, fill=bitgreen!10, draw=bitgreen] (tm)
    at (8,0) {\textbf{Tsetlin}\\\textbf{Machines}\\[2pt]\scriptsize Logic clauses};
  \node[box, fill=bitorange!10, draw=bitorange] (lgn)
    at (0,-5) {\textbf{Logic Gate}\\\textbf{Networks}\\[2pt]\scriptsize Learnable gates};
  \node[box, fill=bitpurple!10, draw=bitpurple] (kan)
    at (8,-5) {\textbf{KAN / LUT}\\\textbf{Compilation}\\[2pt]\scriptsize Train complex,\\[-1pt]\scriptsize deploy as tables};

  % Central node
  \node[draw, circle, minimum size=2cm, fill=lightbg,
        font=\small\sffamily\bfseries, align=center, line width=1pt]
    (center) at (4,-2.5) {Hybrid\\Model};

  % Arrows to center
  \draw[arrow, bitblue]   (bnn.south east)  -- (center);
  \draw[arrow, bitgreen]  (tm.south west)   -- (center);
  \draw[arrow, bitorange] (lgn.north east)  -- (center);
  \draw[arrow, bitpurple] (kan.north west)  -- (center);

  % Labels on arrows -- horizontal, anchored near the boxes
  \node[label, anchor=east] at ($(bnn.south east)+(0,-0.3)$) {binary weights};
  \node[label, anchor=west] at ($(tm.south west)+(0,-0.3)$) {clause logic};
  \node[label, anchor=east] at ($(lgn.north east)+(0,0.3)$) {gate learning};
  \node[label, anchor=west] at ($(kan.north west)+(0,0.3)$) {LUT inference};
\end{tikzpicture}
\caption{The four paradigms of bitwise neural computation and their
  convergence toward a hybrid architecture. Each contributes a different
  capability: binary arithmetic, interpretable logic, learnable
  structure, and efficient deployment.}
\label{fig:paradigm-map}
\end{figure}

\begin{enumerate}[leftmargin=2em]
  \item \textbf{Binary Neural Networks (BNNs)}~\cite{courbariaux2016bnn} constrain weights and
    sometimes activations to just two values ($+1$ and $-1$). This turns
    multiplication into an XNOR gate and addition into bit-counting
    (\cref{sec:bnn}).
    Proven, fast, but accuracy drops on complex tasks.

  \item \textbf{Tsetlin Machines (TMs)}~\cite{granmo2018tsetlin} abandon the neuron metaphor
    entirely. Instead of weighted sums, they learn \emph{if--then} rules
    (propositional clauses) using simple counting automata
    (\cref{sec:tsetlin}). They achieve
    state-of-the-art results on the MNIST handwritten-digit
    dataset~\cite{lecun1998mnist} while remaining fully
    interpretable.

  \item \textbf{Logic Gate Networks (LGNs)}~\cite{petersen2022difflogic} take a middle path: they
    keep the layered network structure but replace neurons with actual
    logic gates (AND, OR, XOR, etc.). Crucially, the \emph{type} of each
    gate is learned during training through a differentiable relaxation
    (\cref{sec:lgn}).

  \item \textbf{Kolmogorov-Arnold Networks (KANs) with LUT compilation}~\cite{liu2024kan}
    flip the script: train with maximally expressive, learnable
    activation functions (B-spline curves on every edge), then
    \emph{compile} those curves into lookup tables (LUTs) for inference.
    The result is a network that runs on table lookups---pure memory
    access, no arithmetic (\cref{sec:kan}).
\end{enumerate}

\subsection{How to Read This Article}

This article is written for engineers, not mathematicians. Each paradigm
is introduced in three stages:

\begin{enumerate}[leftmargin=2em]
  \item \textbf{Intuition}---an analogy or diagram that conveys the core
    idea without any equations.
  \item \textbf{Mechanism}---how it actually works, with diagrams showing
    data flow and simple numerical examples.
  \item \textbf{Formalization}---the mathematical description, for those
    who want to implement it or read the original papers.
\end{enumerate}

You can read at whatever depth suits you. The intuition sections alone
give a solid understanding; the math sections are there when you need
them.

\subsection{Scope and Goal}

Our concrete starting point is handwritten digit recognition on the
MNIST dataset~\cite{lecun1998mnist}---a 28$\times$28 pixel grayscale
image mapped to one of ten digits (0--9). MNIST is deliberately simple:
the goal is not to set accuracy records but to build a \emph{proof of
concept} that demonstrates bitwise computation can work, implemented
entirely in Rust, benchmarked against a standard floating-point baseline.

Beyond the proof of concept, this article also surveys how each paradigm
can scale. \Cref{sec:transformers-llms} provides a self-contained
introduction to transformers and large language models;
\cref{sec:scaling-llms} evaluates whether every transformer component
can be made integer-only or bitwise; and \cref{sec:hybrid}
sketches a hybrid architecture that combines the best ideas from all
four paradigms.
